\section{Fourier Series}

\subsection{Periodic Signal Representation}

\subsubsection{Exponential Fourier series}
\[
x(t)=\sum_{n=-\infty}^{\infty}D_n e^{jn\omega_0 t}, \qquad
D_n=\frac{1}{T_0}\int_{t_0}^{t_0+T_0}x(t)e^{-jn\omega_0 t}\,dt.
\]
The exponential Fourier series expands a periodic signal in orthogonal complex exponentials.

\subsubsection{Trigonometric Fourier series}
\[
x(t)=a_0+\sum_{n=1}^{\infty}\left(a_n\cos n\omega_0 t+b_n\sin n\omega_0 t\right).
\]
The trigonometric series is useful for real periodic signals and separates cosine and sine content.

\subsubsection{Relation between exponential and trigonometric coefficients}
\[
D_0=a_0,\qquad
D_n=\frac{a_n-jb_n}{2},\qquad
D_{-n}=\frac{a_n+jb_n}{2}.
\]
For real signals, \(D_{-n}=D_n^*\). Magnitude and phase follow from the complex coefficients.

\subsubsection{Amplitude-phase Fourier series}
\[
x(t)=C_0+\sum_{n=1}^{\infty}C_n\cos(n\omega_0 t+\theta_n),
\qquad
C_n=2|D_n|,\quad \theta_n=\angle D_n.
\]
The amplitude-phase form describes the spectrum by harmonic amplitudes and phases.

\subsection{Symmetry}

\subsubsection{Fourier series of real even signals}
\[
x(t)=x(-t) \quad \Rightarrow \quad D_n\in\mathbb{R},\qquad b_n=0.
\]
Real even periodic signals contain only cosine terms in the trigonometric form.

\subsubsection{Fourier series of real odd signals}
\[
x(t)=-x(-t) \quad \Rightarrow \quad D_n\in j\mathbb{R},\qquad a_n=0.
\]
Real odd periodic signals contain only sine terms in the trigonometric form.

\subsubsection{Parseval relation for Fourier series}
\[
\frac{1}{T_0}\int_{t_0}^{t_0+T_0}|x(t)|^2\,dt
=\sum_{n=-\infty}^{\infty}|D_n|^2.
\]
Average signal power equals the sum of powers in all Fourier series coefficients.

\subsubsection{Aperiodic time window reconstructed from Fourier series}
\[
x_p(t)=\sum_{k=-\infty}^{\infty}x_a(t-kT_0).
\]
A Fourier series computed from a finite time window represents the periodic extension of that window.
